\section{Introduction}
\label{sec:intro}

\paragraph{Why humor, and why context length.}
Humor is a useful stress test for language-model alignment.  Unlike factual question answering, comedic writing is judged on style, surprise, and topical specificity rather than on verifiable correctness, and good outputs are diverse.  Good comedy is also rooted in context: stand-up comedians do not invent material in a vacuum, they riff on what is happening now and often connect today's events to something that already happened, producing a sort of full-circle callback.  We test whether either ingredient transfers to a fine-tuned language model by varying the prompt context length: a \emph{short} prompt that grounds the model in a current news article only, and a \emph{long} prompt that additionally injects a model-generated summary of related older coverage.  If contextual grounding matters the way it does for human comedians, longer context should produce richer comedic takes; if it doesn't, that is a finding in itself.

\paragraph{Related work.}
RLHF \cite{ouyang2022training,stiennon2020learning} showed that even noisy preference labels can steer large models.  Direct Preference Optimization (DPO) \cite{rafailov2023direct} reformulates the RLHF objective as a closed-form loss over a reference policy, eliminating the explicit reward model.  Kahneman--Tversky Optimization (KTO) \cite{ethayarajh2024kto} generalizes to binary thumbs-up/thumbs-down labels using a prospect-theoretic utility, removing the need for paired preferences.  The closest prior study on humor is \cite{zhang2024humor}, which benchmarks PPO and DPO on multimodal cartoon captioning and reports clear limitations of both; KTO is not evaluated there.  How preference optimization behaves on \emph{text-only} creative generation, and how the prompt context length interacts with the algorithm choice, has received little empirical attention.

\paragraph{Approach.}
We compare DPO and KTO on a comedy-writing task built from Guardian news articles.  We sample two candidate responses per prompt from a frozen Mistral-7B-Instruct-v0.2 \cite{jiang2023mistral}, score each on a three-axis humor rubric using an external judge (Claude Sonnet 4.6), and derive paired (for DPO) and binary (for KTO) datasets from the same scored pool.  Four LoRA-adapted variants \cite{hu2021lora,dettmers2023qlora} are then trained: DPO and KTO under both \emph{short} (headline + lede only) and \emph{long} (plus a model-generated backstory from related older articles) prompts.  Our DPO and KTO losses are implemented from scratch in PyTorch (no Transformer Reinforcement Learning (TRL) or other preference-optimization library), which lets us directly instrument internal quantities such as policy-vs-reference log-probability ratios that high-level training APIs hide.  Evaluation uses an LLM cross-judge on 30 held-out prompts (rank-3 forced choice) plus a parallel blinded ranking by the author.

\paragraph{Contributions.}
(i) Reproducible from-scratch DPO and KTO implementations that match the published formulations and expose internal training metrics for diagnostic plotting.  (ii) Empirical evidence that prompt context length matters as much as the algorithm choice: short-context training improves both algorithms over the base model while long-context training degrades them, under identical hyperparameters.  (iii) A concrete failure mode for KTO: its validation loss falls monotonically while the policy drifts far from the reference, but the bounded loss never registers the drift, so val-loss alone is unsafe as a stopping criterion.  (iv) An algorithm-context interaction: DPO beats KTO head-to-head on short contexts, KTO beats DPO on long contexts.
